\chapter{Préparation, Planification et Architecture Système}

\section{Introduction}
Le succès d'un projet d'intelligence artificielle appliqué à la smart city repose autant sur la rigueur de sa gestion, l'organisation de ses équipes et la fiabilité de son code que sur la perf[...]

\section{Contexte et Motivation du Projet}

\subsection{Contexte du projet et ODD}
Notre projet CIRCU·LI s’inscrit directement dans le contexte mondial des Objectifs de Développement Durable (ODD) définis par les Nations Unies. Il prend place au sein d’un environnement urb[...]

Le projet CIRCU·LI contribue directement et concrètement aux ODD suivants :
\begin{itemize}
    \item \textbf{ODD 9 (Industrie, Innovation et Infrastructure)} : CIRCU·LI introduit des solutions innovantes de gestion des flux urbains. Contrairement aux méthodes archaïques de comptage m[...]
    \item \textbf{ODD 11 (Villes et Communautés durables)} : En améliorant la fluidité du trafic, en réduisant la congestion et en permettant aux urbanistes d'identifier les goulots d'étrangl[...]
    \item \textbf{ODD 13 (Lutte contre le Changement Climatique)} : Les embouteillages sont une source majeure d'émissions de CO2. En optimisant les flux de circulation via des recommandations ba[...]
    \item \textbf{ODD 8 (Travail décent et croissance économique)} : La réduction de la congestion routière diminue les retards professionnels et la perte de productivité. De plus, la réduct[...]
    \item \textbf{ODD 17 (Partenariats pour la réalisation des objectifs)} : Le déploiement de CIRCU·LI repose sur une synergie institutionnelle, favorisant les partenariats stratégiques entre[...]
\end{itemize}

\subsection{Problématique}
La gestion de la mobilité urbaine constitue aujourd’hui l’un des défis les plus complexes des collectivités locales. Avec l’augmentation continue du nombre de véhicules, l’étalement u[...]

Paradoxalement, les municipalités disposent déjà d'infrastructures d'observation (caméras de vidéosurveillance, capteurs passifs). Cependant, ces flux vidéo sont généralement utilisés uni[...]

\subsection{Motivation}
L’évolution des modes de déplacement exige aujourd’hui des solutions plus intelligentes, capables de traiter l'information en temps réel. Notre motivation profonde à développer CIRCU·LI [...]

\section{Conduite du changement et Stratégie (BSC)}

\subsection{Conduite du changement}
L’introduction d’une plateforme d'intelligence artificielle au sein d'une administration municipale n'est pas qu'un défi technique ; c'est un bouleversement organisationnel. La conduite du ch[...]
\begin{enumerate}
    \item \textbf{Identification des acteurs concernés} : Cartographie des parties prenantes (maires, responsables techniques, opérateurs de mobilité, citoyens).
    \item \textbf{Sensibilisation et communication} : Organisation d'ateliers pour démystifier l'IA, mettre en avant la protection de la vie privée (le système compte des véhicules, il n'ident[...]
    \item \textbf{Formation et montée en compétences} : Formation des agents municipaux à la lecture des tableaux de bord interactifs (Business Intelligence) pour en tirer des conclusions urban[...]
    \item \textbf{Mise en œuvre progressive (Proof of Concept)} : Déploiement initial sur un périmètre restreint (ex: un seul grand carrefour stratégique) pour valider l'impact avant la gén�[...]
    \item \textbf{Évaluation et amélioration continue} : Mesure du taux d'adoption de la plateforme par les équipes techniques et recueil de leurs retours d'expérience.
\end{enumerate}

\subsection{Alignement stratégique (Balanced Scorecard - BSC)}
La méthode du Balanced Scorecard (Tableau de bord prospectif) permet de traduire la vision de CIRCU·LI en objectifs tangibles répartis sur quatre axes stratégiques :
\begin{itemize}
    \item \textbf{Axe Financier} : Optimiser le retour sur investissement des infrastructures vidéo existantes. Cible : Éviter l'achat coûteux de capteurs physiques dans la chaussée, réduisan[...]
    \item \textbf{Axe Client (Citoyens et Décideurs)} : Offrir une fluidité de circulation accrue. Cible : Fournir un outil fiable avec un taux de satisfaction des utilisateurs administratifs su[...]
    \item \textbf{Axe Processus Internes} : Automatiser intégralement l'ingestion, le traitement et l'affichage des données. Cible : Garantir une disponibilité (Uptime) de la plateforme supéri[...]
    \item \textbf{Axe Apprentissage et Innovation} : Moderniser l'administration locale. Cible : Diffuser une culture de gestion urbaine "Data-Driven" (basée sur la donnée).
\end{itemize}

\section{Étude comparative et Préparation technologique}

[...]